% from ega4-15-terminal-fr.tex
\oldpage[IV-3]{245}
de $y$ dans $Y$ tel que $f(X)\cap U$ soit fermé dans $U$; alors la restriction $h^{-1}(U)\to U$ de $h$ est une immersion fermée, et la restriction $f^{-1}(U)=g^{-1}(h^{-1}(U))\to h^{-1}(U)$ de $g$ est propre, donc la restriction $f^{-1}(U)\to U$ de $f$ est propre (\textbf{II}, 5.4.2). Pour voir que c) entraîne b), notons d'abord que, pour tout $y\in f(X)$, il existe, en vertu de c), un voisinage ouvert $U$ de $y$ dans $Y$ tel que $f(X)\cap U$ soit fermé dans $U$, donc $f(X)$ est localement fermé, et par suite ouvert dans son adhérence. Comme cette dernière est l'espace sous-jacent de $Z'$, et que $f$ se factorise en $f=j_0\circ g_0$, où $j_0:Z'\to Y$ est l'injection canonique et $g_0$ est un morphisme de $X$ dans $Z'$, le fait que $g_0(X)=Z''$ (qui est ouvert dans $Z'$) entraîne la factorisation de l'énoncé de b); le fait que $g$ est propre résulte alors du caractère local (sur $Z''$) de cette propriété, et de (\textbf{II}, 5.4.3). Il est clair que c) implique d) en vertu de (15.7.5), l'unicité du prolongement d'une $Y'$-section rationnelle résultant de l'hypothèse que $f$ est séparé (\textbf{I}, 7.2.2). Prouvons enfin que d) entraîne c); en premier lieu, il résulte de d), compte tenu de (\textbf{II}, 7.2.3 et 7.2.4), que $f$ est séparé; on peut alors appliquer le critère (15.7.5, c)), qui montre que $f$ est propre en tout point de $f(X)$.
\end{proof}

\begin{remarks}[15.7.7]
\label{IV.15.7.7-fr}
(i) Dans (15.7.4, c)), (15.7.5, c)) et (15.7.6, d)), on peut se restreindre au cas où l'anneau de valuation discrète $A'$ est \emph{complet} et admet un corps résiduel \emph{algébriquement clos}. En effet, dans la démonstration de (15.7.4), si l'on considère un anneau de valuation discrète $A''$ complet dominant $A'$ et ayant un corps résiduel algébriquement clos ($\mathbf{0}_{\mathrm{III}}$, 10.3.1), l'hypothèse c) est alors vérifiée pour $Y''=\operatorname{Spec}(A'')$, $X''=X\times_Y Y''$ et $f''=f_{(Y'')}$; mais $X''=X'\times_{Y'}Y''$, et il résulte alors de la remarque (\textbf{II}, 7.3.9, (i)) que $f'$ est propre; par suite $Y'$, $X'$ et $f'$ vérifient aussi l'hypothèse c) (\textbf{II}, 7.3.8).

(ii) Compte tenu de (\textbf{II}, 7.3.8), la condition c) de (15.7.5) peut se remplacer par la condition que $f'$ est \emph{propre}.
\end{remarks}

\begin{env}[15.7.8]
\label{IV.15.7.8-fr}
On dit qu'un morphisme de préschémas $f:X\to Y$ est \emph{submersif} s'il est surjectif et si la topologie de $Y$ est égale au quotient de la topologie de $X$ par la relation d'équivalence définie par $f$. On dit que $f$ est \emph{universellement submersif} si pour tout changement de base $Y'\to Y$, le morphisme $f'=f_{(Y')}:X\times_Y Y'\to Y'$ est submersif. Tout morphisme surjectif \emph{universellement ouvert}, ou \emph{universellement fermé}, ou \emph{fidèlement plat et quasi-compact} est universellement submersif (2.3.12). Étant donné un morphisme $f:X\to Y$, on dit qu'un sous-ensemble $Z$ de $X$ est \emph{submersif sur $f(Z)$} si la topologie du sous-espace $f(Z)$ de $Y$ est quotient de la topologie du sous-espace $Z$ de $X$ par la relation d'équivalence définie par $f|Z$. On dit que $Z$ est \emph{universellement submersif sur $f(Z)$} si, pour tout changement de base $g:Y'\to Y$, l'image réciproque $Z'=p^{-1}(Z)$ de $Z$ par la projection $p:X\times_Y Y'\to Y'$ est un ensemble submersif sur $f'(Z')=g^{-1}(f(Z))$. On notera que si le morphisme $f$ admet une $Y$-section $h:Y\to X$, tout ensemble $Z$ contenant $h(Y)$ est \emph{universellement submersif sur $Y$}.
\end{env}

\begin{proposition}[15.7.9]
\label{IV.15.7.9-fr}
Soient $Y$ un préschéma localement noethérien, $y$ un point de $Y$, $f:X\to Y$ un morphisme séparé de type fini. Soit $Z$ une partie localement constructible de $X$ ayant les propriétés suivantes :
\begin{enumerate}
  \item[\rm (i)] Pour toute générisation $y'$ de $y$, $Z_{y'}=Z\cap f^{-1}(y')$ est une composante connexe de $f^{-1}(y')$ et est géométriquement connexe sur $\mathbf k(y')$ (4.5.2).
